% BHU PYQ -- Transcribed & Corrected
% Paper: GGRSE-11 / GGRSE11: Man and Environment
% Exam: B.Sc. (Hons) Semester I (NEP) Examination 2024-25
% Ref Code: CECONF/N/2024-25/26 and 27
% Category: Skill Enhancement Course (SEC)
% Department: Department of Geography, Institute of Science / Faculty of Science

\documentclass[11pt,a4paper]{article}
\usepackage[a4paper,top=2.5cm,bottom=2.5cm,left=2.8cm,right=2.8cm,headheight=14pt]{geometry}
\usepackage{amsmath,amssymb}
\usepackage[T1]{fontenc}
\usepackage[utf8]{inputenc}
\usepackage{lmodern,microtype,fancyhdr,enumitem,hyperref,lastpage,parskip}

\hypersetup{
    colorlinks=true,
    urlcolor=black,
    linkcolor=black,
    pdftitle={GGRSE11: Man and Environment -- BHU Sem I 2024-25 (NEP SEC)}
}

\pagestyle{fancy}
\fancyhf{}
\renewcommand{\headrulewidth}{0.4pt}
\renewcommand{\footrulewidth}{0.4pt}
\fancyhead[L]{\small \textit{Banaras Hindu University}}
\fancyhead[R]{\small \textit{GGRSE11: Man and Environment (Sem I 2024-25)}}
\fancyfoot[C]{\small Page \thepage\ of \pageref{LastPage}}

\newlist{parts}{enumerate}{2}
\setlist[parts,1]{label=(\alph*),leftmargin=*,itemsep=4pt,topsep=2pt}

\newcommand{\pts}[1]{\hfill \textit{[#1]}}

\begin{document}

\begin{center}
  {\large\bfseries BANARAS HINDU UNIVERSITY}\\[2pt]
  {\normalsize B.Sc. (Hons) Semester I Examination, 2024-25 (NEP)}\\[6pt]
  \rule{0.8\linewidth}{0.5pt}\\[6pt]
  {\large\bfseries GEOGRAPHY}\\[4pt]
  {\large\bfseries Paper Code: GGRSE-11 : Man and Environment}\\[2pt]
  {\small Ref. No.: CECONF/N/2024-25/26 and 27}\\[4pt]
  {\normalsize Time: 2:30 Hours \hfill Full Marks: 70}
\end{center}

\rule{\linewidth}{0.4pt}

\smallskip

\noindent \textit{Instructions: Attempt five questions including Question 1 which is compulsory. Answer each part of Question 1 in approximately 200 words and of remaining questions in approximately 1000 words. The figures in the right-hand margin indicate marks. Illustrate your answers with suitable maps and diagrams.}\\[2pt]
\noindent \textit{पाँच प्रश्नों के उत्तर दीजिए जिनमें प्रश्न संख्या 1 अनिवार्य है। प्रश्न संख्या 1 के प्रत्येक भाग का उत्तर लगभग 200 शब्दों में तथा शेष प्रश्नों के उत्तर लगभग 1000 शब्दों में दीजिए। उपांत के अंक पूर्णांक के द्योतक हैं। अपने प्रश्नों को उचित मानचित्रों एवं आरेखों द्वारा प्रदर्शित कीजिए।}

\bigskip

\noindent \textbf{Question 1.} Write short notes on the following / निम्नलिखित पर संक्षिप्त टिप्पणियाँ लिखिए:\pts{14}

\begin{parts}
  \item Homo Habilis\\[1pt]
        होमो हैबिलिस \pts{3$\frac{1}{2}$}

  \item Habitat of Eskimos\\[1pt]
        एस्कीमो का निवास स्थान \pts{3$\frac{1}{2}$}

  \item Tundra Biome\\[1pt]
        टुंड्रा बायोम \pts{3$\frac{1}{2}$}

  \item Concept of biodiversity\\[1pt]
        जैवविविधता की अवधारणा \pts{3$\frac{1}{2}$}
\end{parts}

\bigskip
\rule{\linewidth}{0.2pt}
\bigskip

\noindent \textbf{Question 2.} Discuss in detail the bases and classification of human race in the World.\\[1pt]
विश्व में मानव प्रजाति के आधार एवं वर्गीकरण का विस्तृत वर्णन कीजिए। \pts{14}

\medskip\rule{\linewidth}{0.2pt}\medskip

\noindent \textbf{Question 3.} Give a geographical account of Bushman Tribe.\\[1pt]
बुशमैन जनजाति पर एक भौगोलिक लेख लिखिए। \pts{14}

\medskip\rule{\linewidth}{0.2pt}\medskip

\noindent \textbf{Question 4.} Discuss man's interaction with environment and its characteristics with suitable examples.\\[1pt]
मानव एवं पर्यावरण की अन्त:क्रिया एवं उसकी विशेषताओं का सोदाहरण वर्णन कीजिए। \pts{14}

\medskip\rule{\linewidth}{0.2pt}\medskip

\noindent \textbf{Question 5.} Describe the distribution and characteristics of Savannah biomes.\\[1pt]
सवाना बायोम के वितरण एवं विशेषताओं का विवरण प्रस्तुत कीजिए। \pts{14}

\medskip\rule{\linewidth}{0.2pt}\medskip

\noindent \textbf{Question 6.} Elaborate on forms and functions of the ecosystem.\\[1pt]
पारिस्थितिकी तंत्र के स्वरूपों और कार्यों का विस्तार से वर्णन कीजिए। \pts{14}

\medskip\rule{\linewidth}{0.2pt}\medskip

\noindent \textbf{Question 7.} Explain the factors and processes causing the climate change in the World.\\[1pt]
विश्व में जलवायु परिवर्तन के कारकों एवं प्रक्रियाओं की व्याख्या कीजिए। \pts{14}

\vfill
\rule{\linewidth}{0.4pt}

\begin{center}\small
\href{https://bhu-exam.vercel.app/}{Click here to visit our website}\\[4pt]
\href{https://me-aryan.vercel.app/}{Click here to visit our contributor}
\end{center}

\end{document}
